\chapter{Statistical Quality Control}

\section{Quality Improvement and Statistics}

\emph{Quality} means \emph{fitness for use}. It has two parts:
\begin{itemize}
  \item \emph{Quality of design}: the grade, performance and reliability deliberately built into the product.
  \item \emph{Quality of conformance}: how consistently units meet the design. This means steadily reducing variability and defects.
\end{itemize}
Quality improvement is the systematic removal of waste: scrap, rework, inspection, warranty cost and time spent doing things over.

Statistics helps throughout: designed experiments in development (Chapters~7--8), process-capability studies, control charts for monitoring, and life testing for reliability.

\section{Statistical Quality Control}

Statistical quality control (SQC) is the set of statistical and engineering methods for measuring, monitoring, controlling and improving quality.
\begin{itemize}
  \item W.\,A.\ Shewhart (Bell Laboratories) introduced the control chart in 1924.
  \item H.\,F.\ Dodge and H.\,G.\ Romig developed statistical sampling inspection.
  \item W.\,E.\ Deming and J.\,M.\ Juran spread the methods widely after the Second World War.
\end{itemize}

\section{Statistical Process Control}

Quality cannot be inspected into a product. The process must be right the first time. \emph{Statistical process control} (SPC) is the on-line use of statistics to make a process stable and to improve its capability by cutting variability.

The basic problem-solving tools are the \emph{seven tools}:
\begin{enumerate}[label=(\arabic*)]
  \item histogram (or stem-and-leaf plot);
  \item Pareto chart;
  \item cause-and-effect (fishbone) diagram;
  \item defect-concentration diagram;
  \item control chart;
  \item scatter diagram;
  \item check sheet.
\end{enumerate}
These are the technical side of SPC. The other essential part is attitude: a commitment at all levels to continuous improvement by reducing variability. The control chart is the most powerful tool.

\section{Introduction to Control Charts}

\subsection{Basic Principles}

\begin{itemize}
  \item Every process has natural variability from many small, unavoidable \emph{chance causes} (common causes). A process with only chance causes is \emph{in statistical control}.
  \item Other variability comes from badly adjusted machines, operator errors or defective material. It is usually large compared with the background noise. It signals an \emph{assignable cause} (special cause). A process with assignable causes is \emph{out of control}.
  \item A main aim of SPC is to find assignable causes quickly, so corrective action comes before many nonconforming units are made.
\end{itemize}

A \emph{control chart} plots a quality statistic from samples against sample number (or time).
\begin{itemize}
  \item The \emph{center line} (CL) is the average of the statistic when only chance causes act.
  \item The \emph{upper} and \emph{lower control limits} (UCL, LCL) are set so nearly all points fall between them while the process is in control.
  \item A point outside the limits, or a non-random pattern inside them, is evidence of an assignable cause. Investigate.
\end{itemize}

\begin{definition}[General model of a Shewhart control chart]
Let $W$ be a sample statistic measuring a quality characteristic, with mean $\mu_W$ and standard deviation $\sigma_W$. Then
\[ \mathrm{UCL}=\mu_W+k\sigma_W,\qquad\mathrm{CL}=\mu_W,\qquad\mathrm{LCL}=\mu_W-k\sigma_W. \]
$k$ is the distance of the limits from the center line, in standard deviations. The standard choice is $k=3$ (\emph{three-sigma limits}).
\end{definition}

\textbf{A control chart is a repeated hypothesis test.}
\begin{itemize}
  \item For normal $W$, three-sigma limits give $\alpha=0.0027$ per point. That is the chance a point falls outside when the process is in control.
  \item A point inside the limits is like \emph{failing to reject} $H_0$ (process in control). A point outside is like rejecting it.
\end{itemize}

Two types of chart:
\begin{itemize}
  \item \emph{Variables charts} monitor a continuous measurement (a mean chart plus a range or standard-deviation chart).
  \item \emph{Attributes charts} monitor counts of nonconforming units or of nonconformities (defects).
\end{itemize}

\subsection{Design of a Control Chart}

\begin{example}\label{ex:xbardesign}
A process fills containers with nominal mean $\mu=250$ g and known $\sigma=3$ g. Samples of $n=4$ are taken hourly and $\bar{x}$ is plotted.
\begin{itemize}
  \item $\sigma_{\bar{X}}=\sigma/\sqrt n=1.5$.
  \item $\mathrm{UCL}=250+3(1.5)=254.5$ and $\mathrm{LCL}=250-4.5=245.5$.
\end{itemize}
This chart tests $H_0:\mu=250$ at each time point.
\end{example}

For a given multiple of sigma, the limits narrow as $\sqrt n$ grows. The designer must choose sample size and sampling frequency.
\begin{itemize}
  \item Larger samples detect smaller shifts. If only large shifts matter, small samples are enough.
  \item You usually cannot sample both large and often. Current practice favours small, frequent samples.
\end{itemize}

\begin{example}[Piston rings]
Inside diameter has process mean $74$ mm and known $\sigma=0.01$ mm. Samples of $n=5$ are taken hourly.
\begin{itemize}
  \item $\sigma_{\bar{X}}=0.01/\sqrt5=0.0045$.
  \item $\mathrm{UCL}=74+3(0.0045)=74.0135$ and $\mathrm{LCL}=73.9865$ ($\alpha=0.0027$).
\end{itemize}
Choosing these limits is the same as the critical region for repeatedly testing $H_0:\mu=74$ against $H_1:\mu\ne74$ with $\sigma=0.01$ known.
\end{example}

\subsection{Rational Subgroups}

Choose samples (\emph{subgroups}) so that variation \emph{within} a subgroup is only chance variation, and assignable causes show up \emph{between} subgroups. Then the limits bound the chance variation, and assignable causes push points outside.

Two common approaches:
\begin{itemize}
  \item Each subgroup is units made at the same time (a ``snapshot''). This gives the best chance of detecting shifts.
  \item Each subgroup is a random sample of everything made since the last sample. This is used when the aim is to decide whether to accept the product.
\end{itemize}
Rules:
\begin{itemize}
  \item Subgroups must not straddle assignable changes (for example the end of one shift and the start of the next).
  \item Chart the output of several machines or heads \emph{separately}. Do not pool it.
\end{itemize}

\subsection{Analysis of Patterns on Control Charts}

A chart can signal trouble even when all points are inside the limits, if the points follow a non-random pattern: a long \emph{run} (for example all above the center line), a trend, cycles, or many points near the limits. Eight or more points on one side of the center line is very unlikely by chance.

The \emph{Western Electric rules} formalise this. The process is out of control if any of these happens:
\begin{enumerate}[label=\arabic*.]
  \item One point is outside the three-sigma limits.
  \item Two of three consecutive points are beyond a two-sigma limit (same side of the center line).
  \item Four of five consecutive points are at least one sigma from the center line (same side).
  \item Eight consecutive points are on one side of the center line.
\end{enumerate}
\begin{itemize}
  \item The one-, two- and three-sigma lines divide each side of the chart into zones C, B and A.
  \item Rules 2 and 3 apply to one side at a time. A point above the upper two-sigma line followed by one below the lower two-sigma line is not a signal.
  \item These rules add sensitivity but also raise the false-alarm rate. Interpreting patterns needs knowledge of the process.
\end{itemize}

\section{$\bar{X}$ and $R$ or $S$ Control Charts}

The $\bar{X}$ chart controls the average. The range chart ($R$) or standard-deviation chart ($S$) controls the variability.

If $\mu$ and $\sigma$ are known, the $\bar{X}$ chart has
\[ \mathrm{UCL}=\mu+\frac{3\sigma}{\sqrt n},\qquad\mathrm{CL}=\mu,\qquad\mathrm{LCL}=\mu-\frac{3\sigma}{\sqrt n}. \]
When they are unknown, estimate them from $m$ \emph{preliminary samples} of size $n$, taken while the process is believed to be in control. Use at least $20$ to $25$ samples, with $n$ typically $4$, $5$ or $6$. Let $\bar{x}_i$, $r_i$, $s_i$ be the mean, range and standard deviation of sample $i$. Define
\[ \bar{\bar{x}}=\frac1m\sum_{i=1}^m\bar{x}_i,\qquad\bar{r}=\frac1m\sum_{i=1}^mr_i,\qquad\bar{s}=\frac1m\sum_{i=1}^ms_i. \]

\subsubsection*{Estimating $\sigma$ from the range}

For a normal sample, the \emph{relative range} $W=R/\sigma$ has a distribution that depends only on $n$. Its mean is $d_2$ and its standard deviation is $d_3$ (tabulated). So $\Exp[R]=d_2\sigma$ and $\sigma_R=d_3\sigma$. Hence $\hat\sigma=\bar{R}/d_2$ is unbiased, and $\hat\sigma_R=d_3\bar{R}/d_2$.

\begin{theorem}[$\bar{X}$ and $R$ charts]\label{thm:xbarR}
The $\bar{X}$ chart (from $\bar{R}$):
\[ \mathrm{UCL}=\bar{\bar{x}}+A_2\bar{r},\qquad\mathrm{CL}=\bar{\bar{x}},\qquad\mathrm{LCL}=\bar{\bar{x}}-A_2\bar{r},\qquad A_2=\frac3{d_2\sqrt n}. \]
The $R$ chart:
\[ \mathrm{UCL}=D_4\bar{r},\qquad\mathrm{CL}=\bar{r},\qquad\mathrm{LCL}=D_3\bar{r},\qquad D_3=1-\frac{3d_3}{d_2},\quad D_4=1+\frac{3d_3}{d_2}. \]
If $D_3$ (the lower limit) comes out negative, set it to $0$.
\end{theorem}

\subsubsection*{Estimating $\sigma$ from the standard deviation: $S$ chart}

$S$ is a biased estimator: $\Exp[S]=c_4\sigma$, with $c_4$ close to (but below) $1$. Also $\sigma_S=\sigma\sqrt{1-c_4^2}$. With $\bar{s}$ the average sample standard deviation, $\hat\sigma=\bar{s}/c_4$ is unbiased.

\begin{theorem}[$\bar{X}$ and $S$ charts]\label{thm:xbarS}
The $S$ chart:
\[ \mathrm{UCL}=\bar{s}+\frac{3\bar{s}}{c_4}\sqrt{1-c_4^2},\qquad\mathrm{CL}=\bar{s},\qquad\mathrm{LCL}=\bar{s}-\frac{3\bar{s}}{c_4}\sqrt{1-c_4^2}\ (\text{set to }0\text{ if negative}). \]
The $\bar{X}$ chart (from $\bar{S}$):
\[ \mathrm{UCL}=\bar{\bar{x}}+\frac{3\bar{s}}{c_4\sqrt n},\qquad\mathrm{CL}=\bar{\bar{x}},\qquad\mathrm{LCL}=\bar{\bar{x}}-\frac{3\bar{s}}{c_4\sqrt n}. \]
\end{theorem}

\subsubsection*{Factors for control charts}

\begin{center}
\small
\begin{tabular}{c|ccccc|c}
\toprule
$n$ & $A_2$ & $d_2$ & $D_3$ & $D_4$ & $c_4$ & $A=3/\sqrt n$\\
\midrule
2 & 1.880 & 1.128 & 0 & 3.267 & 0.7979 & 2.121\\
3 & 1.023 & 1.693 & 0 & 2.575 & 0.8862 & 1.732\\
4 & 0.729 & 2.059 & 0 & 2.282 & 0.9213 & 1.500\\
5 & 0.577 & 2.326 & 0 & 2.115 & 0.9400 & 1.342\\
6 & 0.483 & 2.534 & 0 & 2.004 & 0.9515 & 1.225\\
7 & 0.419 & 2.704 & 0.076 & 1.924 & 0.9594 & 1.134\\
8 & 0.373 & 2.847 & 0.136 & 1.864 & 0.9650 & 1.061\\
9 & 0.337 & 2.970 & 0.184 & 1.816 & 0.9693 & 1.000\\
10 & 0.308 & 3.078 & 0.223 & 1.777 & 0.9727 & 0.949\\
\bottomrule
\end{tabular}
\end{center}

\begin{example}\label{ex:xbarR}
Ten preliminary samples of size $n=4$.
\begin{itemize}
  \item Means: $50.1,\ 49.8,\ 50.6,\ 50.3,\ 49.5,\ 50.9,\ 50.0,\ 50.4,\ 49.9,\ 50.5$.
  \item Ranges: $3.2,\ 2.8,\ 4.1,\ 3.5,\ 2.9,\ 3.8,\ 3.0,\ 3.3,\ 2.7,\ 3.7$.
  \item $\bar{\bar{x}}=502.0/10=50.20$ and $\bar{r}=33.0/10=3.30$. For $n=4$: $A_2=0.729$, $D_3=0$, $D_4=2.282$.
  \item $\bar{X}$ chart: $\mathrm{CL}=50.20$, $\mathrm{UCL}=50.20+0.729(3.30)=52.61$, $\mathrm{LCL}=50.20-0.729(3.30)=47.79$.
  \item $R$ chart: $\mathrm{CL}=3.30$, $\mathrm{UCL}=2.282(3.30)=7.53$, $\mathrm{LCL}=0$.
\end{itemize}
All means lie between $47.79$ and $52.61$, and all ranges are below $7.53$. The trial limits show the process is in control, so use them to monitor future production. The estimated process standard deviation is $\hat\sigma=\bar{r}/d_2=3.30/2.059=1.60$.
\end{example}

\begin{pitfall}{trial limits and points outside}
Limits from preliminary data are \emph{trial limits}.
\begin{itemize}
  \item Plot the $m$ preliminary means and ranges. Investigate any point outside the limits.
  \item If an assignable cause is found and removed, discard those samples and \emph{recompute} the limits from the rest.
  \item Study the $R$ (or $S$) chart first. If the variability is not stable, the $\bar{X}$ limits (which use $\bar{R}$) are misleading.
  \item Revise limits periodically, and whenever the process improves.
  \item Never put specification limits on a control chart. Specification limits and control limits are unrelated.
\end{itemize}
\end{pitfall}

\begin{example}[Vane opening]
Twenty samples of five parts each (coded data) give $\bar{\bar{x}}=33.3$ and $\bar{r}=5.8$. For $n=5$, $A_2=0.577$.
\begin{itemize}
  \item Trial $\bar{X}$ limits: $33.3\pm0.577(5.8)$, about $29.95$ and $36.65$.
  \item $R$ chart: $\mathrm{CL}=5.8$, $\mathrm{UCL}=D_4\bar{r}$, $\mathrm{LCL}=0$ (since $D_3=0$ for $n=5$).
  \item Samples $6,8,11,19$ are out of control on the $\bar{X}$ chart, and sample $9$ on the $R$ chart (rule 1).
  \item All causes traced to a defective tool in the wax-molding area. Discard these five samples and recompute the limits.
\end{itemize}
Limits from $\bar{s}$ are almost the same as from $\bar{r}$ for such small subgroups. These first $20$ samples are estimation data. The revised limits are then used to judge future production.
\end{example}

\section{Control Charts for Individual Measurements}

Sometimes each sample is a single unit ($n=1$). Every unit may be measured automatically, output may be slow, or only one measurement exists (a batch process).

The \emph{chart for individuals} estimates variability from the \emph{moving range} of two successive observations, $MR_i=\abs{x_i-x_{i-1}}$. Since $d_2=1.128$ for $n=2$, $\hat\sigma=\overline{MR}/1.128$.

\begin{theorem}[Individuals and moving-range charts]\label{thm:individuals}
Individuals chart:
\[ \mathrm{UCL}=\bar{x}+\frac{3\overline{mr}}{1.128},\qquad\mathrm{CL}=\bar{x},\qquad\mathrm{LCL}=\bar{x}-\frac{3\overline{mr}}{1.128}. \]
Moving-range chart ($n=2$, $D_3=0$, $D_4=3.267$): $\mathrm{UCL}=3.267\,\overline{mr}$, $\mathrm{CL}=\overline{mr}$, $\mathrm{LCL}=0$.
\end{theorem}

\begin{example}
Ten concentration readings: $98.2,\ 99.1,\ 97.8,\ 100.4,\ 99.5,\ 98.7,\ 101.0,\ 99.3,\ 98.0,\ 99.9$.
\begin{itemize}
  \item $\bar{x}=99.19$.
  \item The nine moving ranges are $0.9,1.3,2.6,0.9,0.8,2.3,1.7,1.3,1.9$, so $\overline{mr}=13.7/9=1.52$.
  \item $\hat\sigma=1.52/1.128=1.35$.
  \item Individuals chart: $\mathrm{UCL}=99.19+3(1.35)=103.24$, $\mathrm{LCL}=95.14$.
  \item Moving-range chart: $\mathrm{UCL}=3.267(1.52)=4.97$.
\end{itemize}
All points are in control.
\end{example}

Read the individuals chart like an $\bar{X}$ chart, with these cautions:
\begin{itemize}
  \item Moving ranges are correlated, so apparent runs or cycles on the $MR$ chart need not be real.
  \item The observations are assumed uncorrelated, and non-normality matters more.
  \item The chart is slow to detect small shifts. A one-sigma shift takes about $43.9$ points on average.
  \item Narrower limits are \emph{not} the answer: they greatly increase false alarms. Use a CUSUM or EWMA chart for small shifts (Section~\ref{sec:cusum}).
\end{itemize}

\begin{example}[Concentration]
Twenty hourly readings have $\bar{x}=99.1$ and average moving range $\overline{mr}=2.59$. With $D_3=0$, $D_4=3.267$, $d_2=1.128$:
\begin{itemize}
  \item Moving-range chart: $\mathrm{CL}=2.59$, $\mathrm{LCL}=0$, $\mathrm{UCL}=3.267(2.59)=8.46$.
  \item Individuals chart: $99.1\pm3(2.59)/1.128=99.1\pm6.89$.
  \item No point is outside, so the process is in control.
\end{itemize}
\end{example}

\section{Process Capability}

\emph{Process capability} is how a process performs when it is in control. Is its natural variability acceptable relative to the specifications? It is meaningful only for a \emph{stable} process.

Two graphical tools:
\begin{itemize}
  \item The \emph{tolerance chart}: individual measurements plotted with the specification limits LSL and USL. This is the only chart where specification limits belong.
  \item The \emph{histogram}: shows centering and spread.
\end{itemize}

\begin{definition}[Process capability ratio]
With specification limits USL and LSL and process standard deviation $\sigma$,
\[ \mathrm{PCR}=\frac{\mathrm{USL}-\mathrm{LSL}}{6\sigma}. \]
$6\sigma$ is the \emph{width of the process}. The limits $\mu\pm3\sigma$ are the \emph{natural tolerance limits}. The quantity $(1/\mathrm{PCR})\times100$ is the percentage of the specification width the process uses.
\end{definition}

How to read PCR:
\begin{itemize}
  \item $\mathrm{PCR}>1$: the natural limits lie inside the specifications. Few nonconforming units.
  \item $\mathrm{PCR}=1$, centered and normal: $0.27\%$ ($2700$ parts per million) fall outside.
  \item $\mathrm{PCR}<1$: the process is very yield-sensitive. Many nonconforming units.
\end{itemize}
PCR assumes the process is centered at the nominal value, so it measures \emph{potential} capability. For an off-center process, use the \emph{actual} capability:

\begin{definition}[PCR$_k$]
\[ \mathrm{PCR}_k=\min\left[\frac{\mathrm{USL}-\mu}{3\sigma},\ \frac{\mu-\mathrm{LSL}}{3\sigma}\right]. \]
This is a one-sided ratio, using the specification limit nearest the process mean. $\mathrm{PCR}_k=\mathrm{PCR}$ exactly when the process is centered. $\mathrm{PCR}_k<\mathrm{PCR}$ signals off-center operation.
\end{definition}

For a normal process in control, the nonconforming fractions are $\Prob(X<\mathrm{LSL})=\Phi\big((\mathrm{LSL}-\mu)/\sigma\big)$ and $\Prob(X>\mathrm{USL})=1-\Phi\big((\mathrm{USL}-\mu)/\sigma\big)$. These assume normality and control, and use estimated $\mu,\sigma$. Treat them as approximate guidelines.

\begin{example}
For the process of Example~\ref{ex:xbarR}, let the specification be $50\pm6$ ($\mathrm{LSL}=44$, $\mathrm{USL}=56$). With $\hat\sigma=1.603$ and $\bar{\bar{x}}=50.20$:
\begin{itemize}
  \item $\mathrm{PCR}=12/(6\cdot1.603)=1.25$.
  \item $\mathrm{PCR}_k=\min\big[(56-50.2)/(3\cdot1.603),\ (50.2-44)/(3\cdot1.603)\big]=\min(1.21,1.29)=1.21$.
  \item Fallout: $\Prob(X<44)+\Prob(X>56)=\Phi(-3.87)+\big(1-\Phi(3.62)\big)=0.00005+0.00015=0.0002$. About $200$ parts per million.
\end{itemize}
\end{example}

\begin{example}
A current has specifications $100\pm10$ mA. The process has $\mu=107.0$, $\sigma=1.5$.
\begin{itemize}
  \item $\mathrm{PCR}=20/(6\cdot1.5)=2.22$.
  \item $\mathrm{PCR}_k=(110-107)/(3\cdot1.5)=0.67$.
  \item The process is capable but badly off-center. $\Prob(X>110)=\Prob(Z>2)=0.023$: a $2.3\%$ fallout.
\end{itemize}
\end{example}

\emph{Guidelines.}
\begin{itemize}
  \item A common minimum target is $\mathrm{PCR}=1.33$. Use up to $1.66$ for strength, safety or critical characteristics. Some companies require $\mathrm{PCR}_k\ge2.0$.
  \item A process with $\mathrm{PCR}=\mathrm{PCR}_k=2$ is a \emph{six-sigma process}: the nearest specification is six standard deviations away. If its mean shifts by $1.5\sigma$, $\mathrm{PCR}_k$ falls to $1.5$ and fallout is only about $3.4$ parts per million.
  \item In the automotive industry, $C_p$ and $C_{pk}$ mean PCR and PCR$_k$.
\end{itemize}

\section{Attribute Control Charts}

When a unit is just classified as conforming or nonconforming (defective), or defects per unit are counted, use attribute charts. They usually need larger samples than variables charts.

\subsection{$P$ Chart (Control Chart for Proportions)}

Let $D$ be the number of nonconforming units in a random sample of $n$. $D$ is binomial$(n,p)$. The sample fraction nonconforming $\hat{P}=D/n$ has $\mu_{\hat{P}}=p$ and $\sigma^2_{\hat{P}}=p(1-p)/n$. Estimate $p$ by $\bar{p}=\dfrac1{mn}\sum_{i=1}^mD_i$ from $m$ preliminary samples.

\begin{theorem}[$P$ chart]\label{thm:pchart}
\[ \mathrm{UCL}=\bar{p}+3\sqrt{\frac{\bar{p}(1-\bar{p})}n},\qquad\mathrm{CL}=\bar{p},\qquad\mathrm{LCL}=\bar{p}-3\sqrt{\frac{\bar{p}(1-\bar{p})}n}\ (\text{set to }0\text{ if negative}). \]
\end{theorem}

\begin{itemize}
  \item The limits use the normal approximation to the binomial. For small $p$ it is poor; then take limits from binomial probabilities.
  \item The \emph{$NP$ chart} plots the number nonconforming $n\hat{p}$ itself. Multiply everything by $n$: $\mathrm{CL}=n\bar{p}$, limits $n\bar{p}\pm3\sqrt{n\bar{p}(1-\bar{p})}$.
\end{itemize}

\begin{example}
Twenty samples of $n=200$ contain $160$ nonconforming units in total.
\begin{itemize}
  \item $\bar{p}=160/4000=0.04$ and $\sqrt{0.04\cdot0.96/200}=0.01386$.
  \item $\mathrm{UCL}=0.04+0.0416=0.0816$. $\mathrm{LCL}=0.04-0.0416<0$, so set it to $0$.
  \item A sample with more than $16$ nonconforming units ($>0.0816\cdot200$) signals a problem.
\end{itemize}
\end{example}

\subsection{$U$ Chart (Control Chart for Defects per Unit)}

To monitor the number of \emph{defects} (nonconformities) per unit, such as flaws per yard of cloth, model the count as Poisson. If a sample of $n$ inspection units has $C$ defects in total, then $U=C/n$ is the average defects per unit, with mean $\lambda$ and variance $\lambda/n$. With $\bar{u}=\frac1m\sum U_i$:

\begin{theorem}[$U$ chart]\label{thm:uchart}
\[ \mathrm{UCL}=\bar{u}+3\sqrt{\frac{\bar{u}}n},\qquad\mathrm{CL}=\bar{u},\qquad\mathrm{LCL}=\bar{u}-3\sqrt{\frac{\bar{u}}n}\ (\text{set to }0\text{ if negative}). \]
\end{theorem}

The $C$ chart plots the total defects $C$ in a sample: $\mathrm{CL}=n\bar{u}$, limits $n\bar{u}\pm3\sqrt{n\bar{u}}$. These limits use the normal approximation to the Poisson and are poor when $\lambda$ is small.

\begin{example}
Twenty samples of $5$ circuit boards each show $120$ defects in total.
\begin{itemize}
  \item $\bar{u}=120/100=1.2$ defects per board.
  \item $\mathrm{UCL}=1.2+3\sqrt{1.2/5}=2.67$. $\mathrm{LCL}<0$, so set it to $0$.
\end{itemize}
A chart can show the process is in control while the defect level is still too high. $1.2$ defects per board needs process improvement: study the defect types, often with designed experiments.
\end{example}

\begin{example}[$P$ chart and $U$ chart from the textbook]
\emph{$P$ chart.} Twenty samples of $100$ ceramic substrates contain $800$ defectives ($2000$ units in total across samples), so $\bar{p}=0.40$. Limits: $0.40\pm3\sqrt{0.40(0.60)/100}=0.40\pm0.147$, that is, $0.253$ to $0.547$. All samples are in control, but a $40\%$ defective rate is very poor. Analyse the defects by type. An $NP$ chart plots $D$ directly, with limits $n$ times those of the $P$ chart.

\emph{$U$ chart.} Printed circuit boards: five boards inspected per hour, $20$ samples, about $1.6$ defects per board on average. $\bar{u}$ is the average of the $u_i=c_i/n$. Limits: $\bar{u}\pm3\sqrt{\bar{u}/n}$ (a negative $\mathrm{LCL}$ becomes $0$). The process is in control, but $8$ defects per five boards is too many. A $C$ chart plots total defects $c$, with limits $n$ times those of the $U$ chart.
\end{example}

\section{Control Chart Performance}

Moving the limits farther from the center line lowers the Type~I risk (a false alarm: a point outside when no assignable cause is present). But it raises the Type~II risk (failing to signal a real shift).

A measure of performance is the \emph{average run length} (ARL): the average number of points plotted before one signals. If $p$ is the probability that a single point exceeds the limits, the run length is geometric, and
\[ \mathrm{ARL}=\frac1p. \]

\begin{itemize}
  \item \textbf{In control.} For a three-sigma $\bar{X}$ chart $p=0.0027$, so $\mathrm{ARL}_0=1/0.0027=370$. A false alarm comes on average once every $370$ samples.
  \item \textbf{Out of control.} If the mean shifts, $p=1-\beta$ with $\beta=\Prob(\mathrm{LCL}\le\bar{X}\le\mathrm{UCL}\mid\text{new mean})$. If the mean shifts exactly to the upper control limit, $\beta=\Prob(-6\le Z\le0)=0.5$. So $p=0.5$ and $\mathrm{ARL}=2$: detection after two samples on average.
\end{itemize}
The average time to signal is $\mathrm{ARL}\times$ (time between samples). To shorten it, sample more often or increase $n$.

For a shift of $k$ standard deviations of the individual observations, the ARL of a three-sigma $\bar{X}$ chart is:

\begin{center}
\begin{tabular}{c|cc}
\toprule
Shift (multiple of $\sigma$) & ARL, $n=1$ & ARL, $n=4$\\
\midrule
0 & 370.4 & 370.4\\
0.5 & 155.2 & 43.9\\
1.0 & 43.9 & 6.3\\
1.5 & 15.0 & 2.0\\
2.0 & 6.3 & 1.2\\
3.0 & 2.0 & 1.0\\
\bottomrule
\end{tabular}
\end{center}

A Shewhart chart detects large shifts quickly but small shifts slowly. With $n=1$, a $1\sigma$ shift takes $43.9$ points on average. This motivates charts with memory.

\begin{example}[Piston-ring ARL]
The mean shifts to $74.0135$ mm and $n=5$.
\begin{itemize}
  \item The chance that $\bar{X}$ falls between the limits is $0.50$, so $p=0.50$ and $\mathrm{ARL}=1/0.5=2$ samples. On average two hours pass between the shift and its detection.
  \item To detect it faster, sample more often (every half hour gives one hour) or increase $n$. With $n=10$ the limits narrow to $73.9905$ and $74.0095$. The chance of falling between them is about $0.1$, so $p=0.9$ and $\mathrm{ARL}\approx1.11$. About twice as fast.
\end{itemize}
Western Electric rules would lower the in-control ARL below $370$.
\end{example}

\section{Cumulative Sum Control Chart}\label{sec:cusum}

\emph{Cumulative sum} (CUSUM) charts detect \emph{small} shifts better than Shewhart charts, and work well with $n=1$. With $\bar{x}_j$ the mean of sample $j$ and $\mu_0$ the target,
\[ S_i=\sum_{j=1}^i(\bar{x}_j-\mu_0). \]
On target, $S_i$ wanders about $0$. A persistent drift up or down means the mean has shifted. The plain plot has no limits. The older fix is a \emph{V-mask}. The \emph{tabular} CUSUM is preferred.

\begin{definition}[Tabular CUSUM]
With target $\mu_0$, \emph{reference value} $K$, and $S_H(0)=S_L(0)=0$:
\[ S_H(i)=\max\big[0,\ x_i-(\mu_0+K)+S_H(i-1)\big],\qquad S_L(i)=\max\big[0,\ (\mu_0-K)-x_i+S_L(i-1)\big]. \]
\begin{itemize}
  \item $S_H$ and $S_L$ are the \emph{upper} and \emph{lower} one-sided CUSUMs.
  \item They add up deviations beyond $K$ from the target, and reset to $0$ when they would go negative.
  \item The process is out of control if either exceeds the \emph{decision interval} $H$.
\end{itemize}
\end{definition}

After a signal:
\begin{itemize}
  \item If $S_H(i)>H$, estimate the new mean by $\hat\mu=\mu_0+K+S_H(i)/N_H$, where $N_H$ is the number of consecutive periods with $S_H$ nonzero.
  \item If $S_L(i)>H$, use $\hat\mu=\mu_0-K-S_L(i)/N_L$.
  \item Find the cause, then restart the CUSUM at zero.
\end{itemize}

\begin{theorem}[Design of the CUSUM]\label{thm:cusumdesign}
Choose $K$ about half the shift to detect: $K=\delta\sigma_X/2$ for a shift $\delta\sigma_X$. So $K=k\sigma_X$ with $k=\tfrac12$ for a one-sigma shift. Choose $H=h\sigma_X$ with $h=4$ or $5$. Here $\sigma_X$ is the standard deviation of the plotted quantity ($\sigma/\sqrt n$ for sample means). Performance for $k=\tfrac12$:
\begin{center}
\begin{tabular}{c|cc}
\toprule
Shift (multiple of $\sigma_X$) & ARL, $h=4$ & ARL, $h=5$\\
\midrule
0 & 168 & 465\\
0.25 & 74.2 & 139\\
0.50 & 26.6 & 38.0\\
0.75 & 13.3 & 17.0\\
1.00 & 8.38 & 10.4\\
1.50 & 4.75 & 5.75\\
2.00 & 3.34 & 4.01\\
2.50 & 2.62 & 3.11\\
3.00 & 2.19 & 2.57\\
4.00 & 1.71 & 2.01\\
\bottomrule
\end{tabular}
\end{center}
\end{theorem}

\begin{itemize}
  \item For a one-sigma shift the CUSUM needs $8.38$ or $10.4$ samples, against $43.9$ for the Shewhart chart.
  \item The in-control ARL is $168$ for $h=4$ and $465$ for $h=5$. So $h=5$ gives fewer false alarms.
  \item \emph{Standardised CUSUM}: use $y_i=(x_i-\mu_0)/\sigma_X$, with $S_H(i)=\max[0,y_i-k+S_H(i-1)]$ and $S_L(i)=\max[0,-k-y_i+S_L(i-1)]$.
  \item Do \emph{not} apply the Western Electric rules to CUSUM values. Successive values are highly correlated and the rules would give many false alarms.
\end{itemize}

\begin{example}\label{ex:cusum}
Target $\mu_0=50$, individual standard deviation $\sigma=1$. Design for a one-sigma shift with $k=\tfrac12$ and $h=5$: $K=0.5$, $H=5$. Readings: $50.5,\ 49.0,\ 51.8,\ 52.4,\ 50.9,\ 53.1$.
\begin{itemize}
  \item Upper CUSUM $S_H(i)=\max[0,\ x_i-50.5+S_H(i-1)]$: $0,\ 0,\ 1.3,\ 3.2,\ 3.6,\ 6.2$.
  \item Lower CUSUM $S_L(i)=\max[0,\ 49.5-x_i+S_L(i-1)]$: $0,\ 0.5,\ 0,\ 0,\ 0,\ 0$.
  \item At $i=6$, $S_H=6.2>H=5$. Signal an upward shift.
  \item $S_H$ has been nonzero for $N_H=4$ periods ($i=3,\dots,6$). So $\hat\mu=\mu_0+K+S_H/N_H=50+0.5+6.2/4=52.05$.
\end{itemize}
No single reading is far outside three-sigma limits (only $53.1$ would be flagged by $50\pm3$), yet the CUSUM has gathered evidence of a persistent shift.
\end{example}

\begin{example}[Tabular CUSUM from the textbook]
Chemical concentration data: target $\mu_0=99$, $\sigma=2$, $K=1$ ($k=\tfrac12$), $H=10$ ($h=5$). Use $S_H(i)=\max[0,\,x_i-(\mu_0+K)+S_H(i-1)]$ and $S_L(i)=\max[0,\,(\mu_0-K)-x_i+S_L(i-1)]$, both starting at $0$. Neither ever exceeds $H=10$, so the process is in control.
\begin{itemize}
  \item On a signal, find the cause, correct it, and restart both CUSUMs at zero.
  \item The V-mask method is equivalent but less convenient.
  \item Do not apply Western Electric rules to a CUSUM.
  \item The recommended $h=4$ or $5$, $k=\tfrac12$ detects a $1\sigma$ shift in $8.38$ or $10.4$ samples on average.
\end{itemize}
\end{example}

\section{Other SPC Problem-Solving Tools}

The control chart is one of seven basic tools. The others support problem-solving and teamwork.
\begin{itemize}
  \item \textbf{Pareto chart.} Bars of defect types in order of decreasing frequency, with a cumulative-percentage line. It shows the ``vital few, trivial many'' and where effort pays most.
  \item \textbf{Cause-and-effect (fishbone, Ishikawa) diagram.} A team's list of possible causes of a problem, grouped by category (machines, methods, materials, measurement, people, environment). It points to factors worth investigating, often by a designed experiment.
  \item \textbf{Defect-concentration diagram.} A sketch of the product with defect locations marked. Clusters show where, and hence suggest why, defects occur.
  \item \textbf{Histogram.} Shows the shape, center and spread of the data, and how they relate to the specifications.
  \item \textbf{Scatter diagram.} Shows the relationship between two variables (correlation, Chapter~5).
  \item \textbf{Check sheet.} A simple form for collecting data in an organised way at the point of origin.
\end{itemize}
A typical improvement project:
\begin{enumerate}
  \item Rank defect types by Pareto analysis.
  \item Build a fishbone diagram of candidate causes.
  \item Locate a physical source with a defect-concentration diagram.
  \item Use a designed experiment (Chapters~7--8) to find the settings that remove the problem.
\end{enumerate}

\section{Implementing SPC}

Statistical tools alone do not give quality improvement. Successful programmes share these elements:
\begin{enumerate}[label=(\arabic*)]
  \item management leadership and commitment (management is the role model);
  \item a team approach (the basic tools help build teams);
  \item education and training at all levels;
  \item continuous improvement (weekly, quarterly, annual), not a one-time programme;
  \item a way to recognise success.
\end{enumerate}
\begin{itemize}
  \item Introduce control charts early. Processes do not naturally run in control. Using the charts removes assignable causes, reduces variability and stabilises the process. Managing with facts replaces managing by judgement alone.
  \item An \emph{out-of-control-action plan} is essential. A control chart only detects assignable causes. Operators and engineers must find and remove the root cause, not apply a cosmetic fix.
\end{itemize}

\begin{summarybox}
\begin{itemize}
  \item Chance causes are inherent (in control). Assignable causes need action (out of control). A control chart has CL $=\mu_W$ and limits $\mu_W\pm3\sigma_W$ ($\alpha=0.0027$). It is a repeated hypothesis test.
  \item Western Electric rules: one point beyond $3\sigma$; two of three beyond $2\sigma$; four of five beyond $1\sigma$; eight in a row on one side.
  \item $\bar{X}$--$R$: $\bar{\bar{x}}\pm A_2\bar{r}$; $R$: $D_3\bar{r}$, $D_4\bar{r}$; $\hat\sigma=\bar{r}/d_2$. $\bar{X}$--$S$: $\bar{\bar{x}}\pm3\bar{s}/(c_4\sqrt n)$; $S$: $\bar{s}\pm3(\bar{s}/c_4)\sqrt{1-c_4^2}$. Individuals: $\bar{x}\pm3\overline{mr}/1.128$; $MR$: $3.267\overline{mr}$.
  \item Capability: $\mathrm{PCR}=(\mathrm{USL}-\mathrm{LSL})/6\sigma$; $\mathrm{PCR}_k=\min[(\mathrm{USL}-\mu)/3\sigma,(\mu-\mathrm{LSL})/3\sigma]$. Targets $1.33$, $1.66$, $2.0$. Never plot specifications on a control chart.
  \item Attributes: $P$ chart $\bar{p}\pm3\sqrt{\bar{p}(1-\bar{p})/n}$; $U$ chart $\bar{u}\pm3\sqrt{\bar{u}/n}$ (and the $NP$ and $C$ versions).
  \item $\mathrm{ARL}=1/p$. $\mathrm{ARL}_0=370$ for three-sigma limits. CUSUM: $S_H=\max[0,x-(\mu_0+K)+S_H]$, $S_L=\max[0,(\mu_0-K)-x+S_L]$, with $K=k\sigma_X$, $H=h\sigma_X$, $k=\frac12$, $h=4$ or $5$. Much better than Shewhart for small shifts.
\end{itemize}
\end{summarybox}

\section{Exercises}

\begin{exercise}
Twenty-five samples of size $n=5$ have $\sum\bar{x}_i=1250.0$ and $\sum r_i=60.0$. Compute trial control limits for $\bar{X}$ and $R$ charts, and estimate the process standard deviation. If the specifications are $50\pm5$, compute PCR and PCR$_k$.
\end{exercise}

\begin{exercise}
For $n=5$, derive $D_3$ and $D_4$ from $d_2=2.326$ and $d_3=0.864$, and $A_2$ from $d_2$. Show that $B_4=1+3\sqrt{1-c_4^2}/c_4=2.089$ for $n=5$.
\end{exercise}

\begin{exercise}
A $\bar{X}$ chart with three-sigma limits and $n=4$ is used. Compute the probability that a single point plots outside the limits when the process mean shifts by $1.5\sigma$, and the resulting ARL; compare with the table.
\end{exercise}

\begin{exercise}
Twenty-two successive observations on viscosity have $\bar{x}=502.4$ and $\overline{mr}=12.3$. Compute limits for the individuals chart and the moving-range chart, and say whether an observation of $543$ would signal.
\end{exercise}

\begin{exercise}
Twenty-five samples of $n=100$ items contain $175$ defective items in total. Construct a $P$ chart. If the next sample has $13$ defectives is the process in control? What sample size would make the lower control limit positive?
\end{exercise}

\begin{exercise}
A process has USL $=24$, LSL $=12$, $\mu=19$ and $\sigma=1.2$. Compute PCR and PCR$_k$, the fractions nonconforming above USL and below LSL, and comment on centering.
\end{exercise}

\begin{exercise}
Use the tabular CUSUM with $\mu_0=100$, $K=2$ and $H=8$ on the data $101,103,104,102,106,105,108,107,109$. When is a signal given, and what is the estimated new mean?
\end{exercise}

\begin{exercise}
Explain how a Pareto chart and a fishbone diagram would be used in a project to reduce solder defects on circuit boards, and where a control chart and a designed experiment would then enter.
\end{exercise}
